% Dokument pokazowy klasy epi-praca.
%
% Zawiera wszystkie elementy skladu, ktore moga wystapic w pracy licencjackiej:
% strone tytulowa, strony informacyjne, tabele, rysunek, wykres, listing kodu,
% wzor numerowany, cytat blokowy, indeks nazwisk oraz komplet postaci powolan.
%
% Sluzy do sprawdzania klasy i stylu bibliograficznego. Praca powstaje w main.tex.

\documentclass[stronyinfo]{epi-praca}

\addbibresource{bibliografia/literatura.bib}

\autor{Jan Kowalski}
\nralbumu{1234567}
\tytul{Pakiet R do analizy wielokryterialnej jako przykład zastosowania
       oprogramowania badawczego w naukach społecznych}
\tytulen{An R package for multi-criteria analysis as an example of research
         software in the social sciences}
\promotor{dr Marek Deja}
\rok{2027}
\adresstrony{https://jankowalski.github.io/PrzykladR/}
\adresrepozytorium{https://github.com/jankowalski/PrzykladR}
\liczbastron{40}

\abstraktpl{Praca przedstawia projekt, implementację i weryfikację autorskiego
pakietu języka R, który udostępnia badaczom nauk społecznych procedurę
wielokryterialnej analizy decyzyjnej. Celem pracy jest sprawdzenie, czy
algorytm opisany w literaturze metodycznej da się zamknąć w narzędziu
spełniającym wymagania repozytorium CRAN i możliwym do użycia bez przygotowania
programistycznego. W części teoretycznej omówiono aparat matematyczny metody
oraz przesłanki jej wyboru. W części projektowej opisano architekturę pakietu,
procedurę generowania danych testowych o znanej strukturze oraz strategię
testowania obejmującą przypadki analityczne, testy własności i testy odporności.
Działanie narzędzia zilustrowano studium przypadku na otwartym zbiorze danych.
Uzyskane wyniki wskazują, że narzędzie odtwarza rezultaty publikowane w
literaturze i pozostaje stabilne na trzech systemach operacyjnych.}

\abstrakten{The thesis presents the design, implementation and verification of
an original R package that provides social science researchers with a
multi-criteria decision analysis procedure. The aim is to establish whether an
algorithm described in the methodological literature can be enclosed in a tool
that meets CRAN repository requirements and remains usable without programming
background. The theoretical part discusses the mathematical apparatus of the
method and the rationale for its selection. The project part describes the
package architecture, the test data generating procedure with a known structure,
and the testing strategy covering analytical cases, property tests and
robustness tests. The tool is illustrated with a case study on an open dataset.
The results indicate that the tool reproduces published findings and remains
stable across three operating systems.}

\slowaklucze{analiza wielokryterialna, oprogramowanie badawcze, pakiet R,
             odtwarzalność badań, wspomaganie decyzji}
\keywords{decision support, multi-criteria analysis, R package,
          reproducible research, research software}

\begin{document}

\stronatytulowa
\stronyinformacyjne

\tableofcontents

% ===========================================================================
\wprowadzenie

Współczesne badania w naukach społecznych coraz częściej wymagają narzędzi
obliczeniowych, których nie da się zastąpić arkuszem kalkulacyjnym. Metody
opisywane w literaturze metodycznej pozostają jednak dla wielu badaczy
niedostępne, ponieważ nie mają gotowych implementacji \parencite{cisek2002filozoficzne}.
Lukę tę wypełniają pakiety oprogramowania badawczego, budowane i dokumentowane
zgodnie ze standardami odtwarzalności \parencite[s.~15-21]{kowalski2010a}.

Celem pracy jest zaprojektowanie i zweryfikowanie pakietu języka R, który
udostępnia procedurę analizy wielokryterialnej badaczowi bez przygotowania
programistycznego. Praca odpowiada na pytanie, jakie warunki musi spełnić
narzędzie, żeby procedura opisana w artykule metodycznym stała się powtarzalna
i możliwa do zweryfikowania przez osoby trzecie.

Praca składa się z trzech rozdziałów. Rozdział pierwszy przedstawia podstawy
metodyczne. Rozdział drugi opisuje architekturę pakietu. Rozdział trzeci
prezentuje studium przypadku.

% ===========================================================================
\chapter{Podstawy metodyczne}

Metody wielokryterialnego wspomagania decyzji porządkują warianty według wielu
kryteriów jednocześnie. \osoba{Hjørland}{Birger} zwraca uwagę, że wybór aparatu
formalnego przesądza o zakresie wniosków, jakie wolno wyciągnąć z analizy
\parencite{hjorland1998theory}. Podobne stanowisko zajmuje
\textcite{wozniak1997kognitywizm}.

\section{Aparat formalny}

Znormalizowaną wartość wariantu wyznacza się jako iloraz odległości od
rozwiązania anty-idealnego i sumy obu odległości:

\begin{equation}\label{eq:cc}
  CC_i = \frac{d_i^{-}}{d_i^{+} + d_i^{-}},
  \qquad
  d_i^{\pm} = \sqrt{\sum_{j=1}^{n} \bigl(v_{ij} - v_j^{\pm}\bigr)^2}.
\end{equation}

We wzorze \eqref{eq:cc} symbol $v_{ij}$ oznacza znormalizowaną i zważoną ocenę
wariantu $i$ względem kryterium $j$, natomiast $v_j^{+}$ i $v_j^{-}$ to wartości
odpowiednio najlepsza i najgorsza w obrębie kryterium. Wskaźnik $CC_i$ przyjmuje
wartości z przedziału $[0, 1]$; im bliżej jedności, tym wariant lepszy.

\section{Wyznaczanie wag}

Wagi obiektywne wyprowadza się z rozproszenia ocen w obrębie kryterium. Jak
ujmuje to jedno z opracowań podstawowych:

\begin{quote}
Kryterium, w którym oceny wariantów są do siebie zbliżone, niesie mało
informacji różnicującej i powinno otrzymać wagę niższą niż kryterium
o wysokim rozproszeniu ocen.
\end{quote}

Podejście to opisano szerzej w pracach zbiorowych poświęconych metodyce badań
\parencite{kowalski1998red} oraz w opracowaniach przeglądowych
\parencite{nowak1999wielu,nowakp1999}.

% ===========================================================================
\chapter{Implementacja i architektura pakietu}

Pakiet zaimplementowano w języku R zgodnie ze standardami repozytorium CRAN.
Logika obliczeniowa znajduje się w katalogu \plik{R/}, dane w \plik{data/},
a dokumentacja w \plik{man/} oraz \plik{vignettes/}.

\section{Struktura i przetwarzanie danych}

Funkcja \fun{przygotuj\_dane} przyjmuje surową ramkę danych oraz definicję
modelu i zwraca macierz decyzyjną. Argument \argument{skala} określa zakres
normalizacji, a argument \argument{braki} sposób obsługi wartości brakujących.
Wizualizacje powstają z użyciem pakietu \pkg{ggplot2}.

Warto podkreślić, że \wyroznienie{walidacja wejścia poprzedza każde obliczenie} --
jest to \termin{precondition} w rozumieniu programowania kontraktowego.

\begin{tabelaepi}{Moduły pakietu i ich odpowiedzialność}{tab:moduly}
\begin{tabular}{@{}lll@{}}
\toprule
Plik & Odpowiedzialność & Funkcje eksportowane \\
\midrule
\plik{przygotowanie\_danych.R} & walidacja i normalizacja & \fun{przygotuj\_dane} \\
\plik{algorytm.R}              & silnik obliczeniowy      & \fun{oblicz\_ranking} \\
\plik{klasy\_s3.R}             & obiekt wyniku            & \fun{print}, \fun{summary} \\
\plik{wizualizacja.R}          & mapy decyzyjne           & \fun{plot} \\
\bottomrule
\end{tabular}
\zrodlo{oprac. własne}
\end{tabelaepi}

Pełny przebieg analizy przedstawia listing~\ref{lst:workflow}.

\begin{lstlisting}[caption={Ścieżka analityczna od danych surowych do rankingu},label={lst:workflow}]
library(PrzykladR)

# Definicja modelu w skladni zblizonej do lavaan
model <- "Koszt =~ koszt_surowce + koszt_praca
          Jakosc =~ jakosc_trwalosc + jakosc_ux"

dane <- przygotuj_dane(mcda_dane_surowe, model, skala = c(1, 9))
wynik <- oblicz_ranking(dane, wagi = "entropia")

summary(wynik)
plot(wynik)
\end{lstlisting}

\section{Prezentacja wyników}

\begin{rysunekepi}{Mapa decyzyjna wariantów w przestrzeni odległości}{rys:mapa}
\includegraphics[width=0.78\textwidth]{przyklad-wykres.png}
\zrodlo{oprac. własne}
\end{rysunekepi}

Rysunek~\ref{rys:mapa} pokazuje położenie wariantów względem obu punktów
odniesienia. Wariant oznaczony literą A znajduje się w ćwiartce najkorzystniejszej.

\begin{wykresepi}{Udział kryteriów w wagach obiektywnych}{wyk:wagi}
\begin{tabular}{@{}lr@{}}
\toprule
Kryterium & Waga \\
\midrule
Koszt   & 0,31 \\
Jakość  & 0,28 \\
Dostawa & 0,24 \\
Rozwój  & 0,17 \\
\bottomrule
\end{tabular}
\zrodlo{oprac. własne na podstawie \parencite{wozniak1997kognitywizm}}
\end{wykresepi}

% ===========================================================================
\chapter{Studium przypadku}

Weryfikację przeprowadzono na zbiorze wygenerowanym własną procedurą
o znanej strukturze zależności. Procedurę opisano w normie dotyczącej
dokumentowania źródeł \parencite{pniso690} oraz w dokumentacji środowiska
\parencite{rcore2026}.

Zastosowane hasła terminologiczne opierają się na wydawnictwach informacyjnych
\parencite{apoteoza1996,kolodziejska1976propaganda} oraz na opracowaniach
dostępnych w wersji elektronicznej \parencite{cisek2002szara,kowalski2009zmiany}.
Materiały instytucjonalne przywołano za serwisem biblioteki
\parencite{bj2021zasoby} i wpisem badaczki \parencite{cisek2019informacja}.

Pozostałe pozycje sprawdzające postać opisu to praca zbiorowa
\parencite{zielinski1998swiat}, rozdział w tomie \parencite{nowak2004srodowisko},
artykuł prasowy \parencite{nowak2000jutro}, wydawnictwo na nośniku
\parencite{kopalinski1998slownik}, pozycja bez autora
\parencite{elektroniczne2008}, druga praca tego samego autora z tego samego roku
\parencite{kowalski2010b} oraz książka dwóch autorów
\parencite{nowak2009postrzeganie}.

% ===========================================================================
\podsumowanie

Zbudowane narzędzie odtwarza procedurę opisaną w literaturze metodycznej
i przechodzi sprawdzenie zgodności z wymaganiami repozytorium na trzech
systemach operacyjnych. Dalszy rozwój obejmuje rozszerzenie zestawu metod
wyznaczania wag oraz udostępnienie interfejsu przeglądarkowego.

% ===========================================================================
\wykazzrodel

Wykaz obejmuje narzędzia programistyczne wykorzystane w projekcie.

\begin{description}
  \item[\pkg{ggplot2}] Wickham, H. (2016). \termin{ggplot2: Elegant Graphics for
    Data Analysis}. Nowy Jork: Springer.
  \item[\pkg{testthat}] Wickham, H. (2011). testthat: Get Started with Testing.
    \termin{The R Journal}, vol. 3, pp. 5-10.
  \item[\pkg{PrzykladR}] Kowalski, J. (2027). \termin{PrzykladR: Wielokryterialna
    analiza decyzyjna}. \url{https://github.com/jankowalski/PrzykladR}.
\end{description}

% ===========================================================================
\bibliografia

% ===========================================================================
\spisilustracji

% ===========================================================================
\indeksnazwisk

% ===========================================================================
\aneksy

\aneks{Metadane oprogramowania}

\begin{tabular}{@{}ll@{}}
\toprule
Nr & Opis metadanych \\
\midrule
C1 & Nazwa oprogramowania: PrzykladR \\
C2 & Wersja: 1.0.0 \\
C3 & Licencja: GPL-3 \\
C4 & Repozytorium: \url{https://github.com/jankowalski/PrzykladR} \\
C5 & Język programowania: R \\
C6 & Zależności: ggplot2, stats \\
\bottomrule
\end{tabular}

\aneks{Nota wymagana w serwisie WWW}

Autorem niniejszego serwisu jest Jan Kowalski. Serwis stanowi integralną część
pracy licencjackiej (kierunek: elektroniczne przetwarzanie informacji),
przygotowanej pod kierunkiem dra Marka Dei na Wydziale Zarządzania
i Komunikacji Społecznej Uniwersytetu Jagiellońskiego.

\end{document}
